\chapter{2010年陕西卷（理）}

\section{单选题}

\begin{problem}
集合 \(A = \{x \mid -1 \leqslant x \leqslant 2\}\)，\(B = \{x \mid x < 1\}\)，则 \(A \cap (\complement_{\R} B) =\)（\quad）

\choices
    {\( \{x \mid x > 1\} \)}
    {\(\{x \mid x \geqslant 1\}\)}
    {\(\{x\mid 1 < x\leqslant 2\}\)}
    {\(\{x\mid 1\leqslant x\leqslant 2\}\)}
\end{problem}

\begin{problem}
复数 \(z = \frac{1}{1 + \i}\) 在复平面上对应的点位于（\quad）

\choices
    {第一象限}
    {第二象}
    {第三象限}
    {第四象限}
\end{problem}

\begin{problem}
对于函数 \( f(x) = 2\sin x\cos x \)，下列选项中正确的是（\quad）

\choices
    {对于 \(f(x)\) 在 \(\left(\frac{\pi}{4}, \frac{\pi}{2}\right)\) 上是递增的}
    {\(f(x)\) 的图象关于原点对称}
    {\(f(x)\) 的最小正周期为 \(2\pi\)}
    {\(f(x)\) 的最大值为2}
\end{problem}

\begin{problem}
\(\left(x + \frac{y}{2}\right)^3 (x\in \R)\) 展开式中 \(x^{3}\) 的系数为10，则实数a等于（\quad）

\choices
    {\(-1\)}
    {\( \frac{1}{2} \)}
    {1}
    {2}
\end{problem}

\begin{problem}
已知函数 \( f(x) = \begin{cases} 2^x + 1, & x < 1, \\ x^2 + ax, & x \geqslant 1, \end{cases} \) 若 \( f(f(0)) = 4n \)，则实数 \( a \) 等于（\quad）

\choices
    {\(\frac{1}{2}\)}
    {\(\frac{4}{5}\)}
    {2}
    {9}
\end{problem}

\begin{problem}
如图所示是求样本 \(x_{1}, x_{2}, \cdots, x_{10}\) 平均数 \(\overline{x}\) 的程序框图，图中空白框中应填入的内容为（\quad）

\choices
    {\(S = S + x_{n}\)}
    {\(S = S + \frac{x_n}{n}\)}
    {\(S = S + n\)}
    {\(S = S + \frac{1}{n}\)}
\end{problem}

\begin{problem}
若某空间几何体的三视图如图所示，则该几何体的体积是（\quad）

俯视图（\quad）

\choices
    {\(\frac{1}{3}\)}
    {\(\frac{2}{3}\)}
    {1}
    {2}
\end{problem}

\begin{problem}
已知抛物线 \( y^{2}=2px \) (p>0) 的准线与圆 \( x^{2}+y^{2}-6x-7=0 \) 相切，则 p 的值为（\quad）

\choices
    {\( \frac{1}{2} \)}
    {1}
    {2}
    {4}
\end{problem}

\begin{problem}
对于数列 \(\{a_{n}\}, "a_{n+1} > |a_{n}| (n = 1, 2, \cdots)\) 是“\([a_{n}\}\) 为递增数列”的（\quad）

\choices
    {必要不充分条件}
    {充分不必要条件}
    {充要条件}
    {既不充分也不必要条件}
\end{problem}

\begin{problem}
某学校要开拓学生代表大会，规定各班每 10 人推选一名代表，当各班人数除以 10 的余数大于 6 时得选出一名代表. 那么，各班可推选代表人数 \( n \) 与该班人数 \( x \) 之间的函数关系用取整函数 \( y = |x| \) ( \( |x| \) 表示不大于 \( x \) 的最大整数) 可以表示为（\quad）

\choices
    {\( y = \left[\frac{x}{10}\right] \)}
    {\( y = \left[\frac{x + 3}{10}\right] \)}
    {\( y = \left[\frac{x + 4}{10}\right] \)}
    {\( y = \left[\frac{x + 5}{10}\right] \)}
\end{problem}

\section{填空题}

\begin{problem}
已知向量 \( a = (2, -1) \) ，\( b = (-1, m) \) ，\( c = (-1, 2) \) . 若 \( (a + b) // c \) ，则 \( m = \)  \fillinblank{}.
\end{problem}

\begin{problem}
观察下列等式: \(1^{1} + 2^{1} = 3^{2}, 1^{3} + 2^{3} + 3^{3} = 6^{2}, 1^{3} + 2^{3} + 3^{3} + 4^{3} = 10^{2}\)， \(\cdots\)，根据上述规律，第五个等式为 \fillinblank{} \fillinblank{}.
\end{problem}

\begin{problem}
从如图所示的长方形区域内任取一个点 \(M(x, y)\)，则点 \(M\) 取自阴影部分的概率为 \fillinblank{} \fillinblank{}.
\end{problem}

\begin{problem}
铁矿石 \(A\) 和 \(B\) 的含铁率 \(\alpha_{1}\) 冶炼每万吨铁矿石的 \(\mathrm{CO}_{2}\) 排放量 \(b\) 及每万吨铁矿石的价格 \(c\) 如下表： b(万吨) c(百万元) A 50\% 1 3 B 70\% 0.5 6 某冶炼厂至少要生产1.9(万吨)铁，若要求 \( CO_{2} \) 的排放量不超过2(万吨)，则购买铁矿石的最少费用为(百万元) \fillinblank{}.
\end{problem}

\begin{problem}
三选一. 【A】不等式 \(|x + 3| - |x - 2|\geqslant 3\) 的解集为 【B】如图，已知 Rt\(\triangle\)ABC 的两条直角边 AC, BC 的长分别为 3 cm, 4 cm，以 AC 为直径的圆与 AB 交于点 D，则 \( \frac{BD}{DA}= \) .

【C】已知圆 C 的参数方程为 \( \left\{\begin{aligned}x&=\cos\alpha,\\ y&=1+\sin\alpha,\end{aligned}\right. \) ( \( \alpha \) 为参数)，以原点为极点，x 轴正半轴为极轴建立极坐标系. 直线 l 的极坐标为 \( p\sin\theta=1 \) ，则直线 l 与圆 C 的交点的直角坐标为 \fillinblank{} \fillinblank{}.
\end{problem}

\section{解答题}

\begin{problem}
已知 \(\{a_{n}\}\) 是公差不为零的等差数列，\(a_{1} = 1\) 且 \(a_{1}, a_{3}, a_{9}\) 成等比数列. (1) 求数列 \( \{a_{n}\} \) 的通项; (2) 求数列 \( \{2^{n}\} \) 的前 n 项和 \( S_{n} \) .
\end{problem}

\begin{problem}
如图，A, B 是海面上位于东西方向相距 \( 5(3 + \sqrt{3}) \) 海里的两个观测点. 现位于 A 点北偏东 \( 45^{\circ} \) ，B 点北偏西 \( 60^{\circ} \) 的 D 点有一艘轮船发出求救信号，位于 B 点南偏西 \( 60^{\circ} \) 且与 B 点相距 \( 20\sqrt{3} \) 海里的 C 点的数据超过靠前往营救，其航行速度为 30 海里/小时，该数据抛到边 D 点需要多长时间？
\end{problem}

\begin{problem}
如图，在四棱锥 \(P - ABCD\) 中，底面 \(ABCD\) 是矩形，\(PA \perp\) 平面 \(ABCD\)，\(AP = AB = 2\)，\(BC = 2\sqrt{2}\)，\(E\)，\(F\) 分别是 \(AD\)，\(PC\) 的中点. (1) 证明: \(PC \perp\) 平面 \(BEF\); (2) 求平面 BEF 与平面 BAP 夹角的大小.
\end{problem}

\begin{problem}
如图，椭圆 \(C: \frac{x^2}{a^2} + \frac{y^2}{b^2} = 1\) 的顶点为 \(A_1, A_2, B_1, B_2\)，焦点为 \(F_1, F_2\)，\(|A_1B_1| = \sqrt{7}, s_{2D} = s_1s_2, b_1 = 2s_2s_1, s_1s_2 = 2s_2s_1, s_1s_2 = s_1s_2\). (1) 求椭圆 C 的方程. (2) 设 n 是过原点的直线，l 是与 n 垂直相交于 P 点、与椭圆相交于 A, B 两点的直线，\( \overrightarrow{OP} = 1 \) . 是否存在上述直线 l 使 \( \overrightarrow{AP} \cdot \overrightarrow{PB} = 1 \) 成立\(\perp\) 若存在，求出直线 l 的方程; 若不存在，请说明理由.
\end{problem}

\begin{problem}
已知函数 \( f(x) = \sqrt{x} \)，\( g(x) = a \ln x \)，\( a \in \R \). (1) 若曲线 \( y = f(x) \) 与曲线 \( y = g(x) \) 相交，且在交点处有共同的切线，求 \( a \) 的值和该切线方程； (2) 设函数 \( h(x) = f(x) - g(x) \)，当 \( h(x) \) 存在最小值时，求其最小值 \( \varphi(a) \) 的解析式； (3) 对 (2) 中的 \(\varphi(a)\) 和任意的 \(a > 0, b > 0\). 证明: \(\varphi'\left(\frac{a + b}{2}\right) \leqslant \frac{\varphi'(a) + \varphi'(b)}{2} \leqslant \varphi'\left(\frac{2ab}{a + b}\right)\)
\end{problem}

\begin{problem}
为了解学生身高情况，某校以 \(10\%\) 的比例对全校 700 名学生按性别进行分层抽样调查，测得身高情况的统计图如下: (1) 估计该校男生的人数; (2) 估计该校学生身高在 \(170 \sim 185 \mathrm{~cm}\) 之间的概率; (3) 从样本中身高在 \(165 \sim 180 \mathrm{~cm}\) 之间的女生中任选 2 人，求至少有 1 人身高在 \(170 \sim 180 \mathrm{~cm}\) 之间的概率.
\end{problem}

